\section{Preliminaries}
\label{sec:pre}

\paragraph{Notations}
Bold font $\vecx$ denotes a vector in $\RR^d$ and $\|\vecx\|$ denotes its Euclidean norm. We define $B_d(\vecx,r)=\{\vecy\in\RR^d:\|\vecx-\vecy\|\le r\}$ and sometimes drop the subscript $d$ when the context is clear. We use $[n]$ as an abbreviation for $\{1,2,\ldots,n\}$. We adopt the standard big-$O$ notation, and $f\lesssim g$ denotes $f=O(g)$. $\calP(S)$ denotes the set of all distributions over a measurable set $S$.


\paragraph{Stochastic Optimization}
Given a function $F:\RR^d\to\RR$, $F$ is $G$-Lipschitz if $|F(\vecx)-F(\vecy)| \le G\|\vecx-\vecy\|, \forall \vecx,\vecy$. Equivalently, when $F$ is differentiable, $F$ is $G$-Lipschitz if $\|\nabla F(\vecx)\|\le G, \forall \vecx$. $F$ is $H$-smooth if $F$ is differentiable and $\nabla F$ is $H$-Lipschitz; $F$ is $\rho$-second-order-smooth if $F$ is twice differentiable and $\nabla^2 F$ is $\rho$-Lipschitz.

\begin{assumption}
\label{as:optimization}
We assume that our objective function $F:\RR^d\to \RR$ is differentiable and $G$ Lipschitz, and given an initial point $\vecx_0$, $F(\vecx_0)-\inf F(\vecx) \le F^*$ for some known $F^*$. We also assume the stochastic gradient satisfies $\E[\nabla f(\vecx,z)\,|\,\vecx]=\nabla F(\vecx), \E\|\nabla F(\vecx)-\nabla f(\vecx,z)\|^2 \le \sigma^2$ for all $\vecx,z$. Finally, we assume that $F$ is \emph{well-behaved} in the sense of \cite{cutkosky2023optimal}: for any points $\vecx$ and $\vecy$, $F(\vecx)-F(\vecy) = \int_0^1 \langle \nabla F(\vecx+t(\vecy-\vecx)), \vecy-\vecx\rangle \ dt$.
% Following the convention in literature, we assume $F$ is $G$-Lipschitz, $F(\vecx_0)-\inf F(\vecx) \le F^*$ and $\E[\nabla f(\vecx,z)\,|\,\vecx]=\nabla F(\vecx), \E\|\nabla F(\vecx)-\nabla f(\vecx,z)\|^2 \le \sigma^2$ for all $\vecx,z$.
\end{assumption}


\paragraph{Online Learning}
An online convex optimization (OCO) algorithm is an iterative algorithm that uses the following procedure: in each iteration $n$, the algorithm plays an action $\Delta_n$ and then receives a convex  loss function $\ell_n$ The goal is to minimize the regret w.r.t. some comparator $\vecu$, defined as
\begin{equation*}
    \regret_n(\vecu) := \littlesum_{t=1}^n \ell_t(\Delta_t) - \ell_t(\vecu).
\end{equation*}
The most basic OCO algorithm is online gradient  descent: $\Delta_{n+1} = \Delta_n - \eta \nabla \ell_n(\Delta_n)$, which guarantees $\regret_N(\vecu) =O(\sqrt{N})$ for appropriate $\eta$. Notably, in OCO we make \emph{no assumptions} about the dynamics of $\ell_n$. They need not be stochastic, and could even be adversarially generated. We will be making use of algorithms that obtain \emph{anytime} regret bounds. That is, for all $n$ and any sequence of $\vecu_1,\vecu_2,\dots$, it is possible to bound $\regret_n(\vecu_n)$ by some appropriate quantities (that may be function of $n$). This is no great burden: almost all online convex optimization algorithms have anytime regret bounds.
% Note that formulation above is technically a simplified description of OCO sometimes called ``online \emph{linear} optimization'' or OLO. 
For readers interested in more details, please refer to \cite{cesa2006prediction, hazan2019introduction, orabona2019modern}.


% \paragraph{Momentum Algorithms}
% We recall the update formula of two popular momentum algorithms, SGDM and Adam. Given stochastic gradient $\vecg_t=\nabla f(\vecx_t,z_t)$, SGDM updates $\m_{t+1} = \beta\m_t + (1-\beta)\vecg_t$ and $\vecx_{t+1} = \vecx_t - \eta\m_{t+1}$.\footnote{The original update in \cite{sutskever2013importance} is $\vecv_{t+1} = \beta\vecv_t - \mu\vecg_t$ and $\vecx_{t+1}=\vecx_t+\vecv_{t+1}$. However, substituting $\mu=(1-\beta)\eta$ and $\vecv_t = \frac{-\mu}{1-\beta}\m_t$ recovers the same update.} Here $\beta$ denotes the momentum constant and $\eta$ denotes the learning rate.

% Similarly, Adam updates $\m_{t+1} = \beta_1\m_t + (1-\beta_1)\vecg_t$, $\vecv_{t+1} = \beta_2\vecv_t + (1-\beta_2)\vecg_t^2$ and $\vecx_{t+1} = \vecx_t - \eta \frac{\m_{t+1}/(1-\beta_1^t)}{\sqrt{\vecv_{t+1}/(1-\beta_2^t)}}$. Here $\vecv_t$ is updated coordinate-wise.


\subsection{Non-smooth Optimization}
% \paragraph{Convergence Criterion}

Suppose $F$ is differentiable. $\vecx$ is an $\epsilon$-stationary point of $F$ if $\|\nabla F(\vecx)\|\le \epsilon$. This definition is the standard criterion for smooth non-convex optimization.
For non-smooth non-convex optimization, the standard criterion is the following: $\vecx$ is an $(\delta,\epsilon)$-Goldstein stationary point of $F$ if there exists $S\subset \RR^d$ and $P\in\calP(S)$ such that $\vecy\sim P$ satisfies $\E[\vecy]=\vecx$, $\|\vecy-\vecx\|\le\delta$ almost surely, and $\|\E[\nabla F(\vecy)]\|\le\epsilon$.\footnote{The original definition of $(\delta,\epsilon)$ Goldstein stationary point proposed by \cite{zhang2020complexity} does not require the condition $\E[\vecy]=\vecx$. However, as shown in \cite{cutkosky2023optimal}, this condition allows us to reduce a Goldstein stationary point to an $\epsilon$-stationary point when the loss is second-order smooth. Hence we also keep this condition.} Next, we formally define $(c,\epsilon)$-stationary point, our proposed new criterion for non-smooth optimization.

\begin{definition}
\label{def:Exp-O2NC:stationary-point}
Suppose $F:\RR^d\to\RR$ is differentiable, $\vecx$ is a $(c,\epsilon)$-stationary point of $F$ if $\|\nabla F(\vecx)\|_c \le \epsilon$, where
\begin{equation*}
    \|\nabla F(\vecx)\|_c = \inf_{\substack{S\subset \RR^d \\ \vecy\sim P\in\calP(S) \\ \E[\vecy]=\vecx}} \| \E[\nabla F(\vecy)]\| + c \cdot \Ex\|\vecy-\vecx\|^2.
\end{equation*}
\end{definition}

In other words, if $\vecx$ is a $(c,\epsilon)$-stationary point, then there exists $S\subset \RR^d, P\in\calP(S)$ such that $\vecy\sim P$ satisfies $\E[\vecy]=\vecx$, $\E\|\vecy-\vecx\|^2 \le \epsilon/c$, and $\|\E[\nabla F(\vecy)]\|\le \epsilon$. To see how this definition is related to the previous $(\epsilon,\delta)$-Goldstein stationary point definition, consider the case when $c=\epsilon/\delta^2$. Then this new definition of $(c,\epsilon)$-stationary point is almost identical to $(\delta,\epsilon)$-Goldstein stationary point, except that it relaxes the constraint from $\|\vecy-\vecx\|\le\delta$ to $\E\|\vecy-\vecx\|^2\le\delta^2$. 

% ELABORATE why this is good definition

To further motivate this definition, we demonstrate that $(c,\epsilon)$-stationary point retains desirable properties of Goldstein stationary points. Firstly, the following result shows that, similar to Goldstein stationary points, $(c,\epsilon)$-stationary points can also be reduced to first-order stationary points with proper choices of $c$ when the objective is smooth or second-order smooth. 

\begin{restatable}{lemma}{CriterionReduction}
\label{lem:Exp-O2NC:criterion-reduction}
% Assume $\|\nabla F(\vecx)\|_c\le \epsilon$ for any $c=\epsilon\delta^{-2}$ and $\delta>0$. Then $\|\nabla F(\vecx)\| \le \epsilon$ if $F$ is $H$-smooth and $\delta=\eps/2H$ or if $F$ is $\rho$-second-order-smooth and $\delta=\sqrt{2\epsilon}/\sqrt{\rho}$.
Suppose $F$ is $H$-smooth. 
If $\|\nabla F(\vecx)\|_c\le\epsilon$ where $c=H^2\epsilon^{-1}$, then $\|\nabla F(\vecx)\|\le 2\epsilon$. \\
% If $\|\nabla F(\vecx)\|_c\le\epsilon$ where $c=\epsilon\delta^{-2}, \delta=\epsilon/2H$, then $\|\nabla F(\vecx)\|\le 2\epsilon$. \\
Suppose $F$ is $\rho$-second-order-smooth. 
If $\|\nabla F(\vecx)\|_c\le\epsilon$ where $c=\rho/2$, then $\|\nabla F(\vecx)\|\le 2\epsilon$.
% If $\|\nabla F(\vecx)\|_c\le\epsilon$ where $c=\epsilon\delta^{-2}, \delta=\sqrt{2\epsilon}/\sqrt{\rho}$, then $\|\nabla F(\vecx)\|\le 2\epsilon$.
\end{restatable}

As an immediate implication, suppose an algorithm achieves $O(c^{1/2}\epsilon^{-7/2})$ rate for finding a $(c,\epsilon)$-stationary point. Then Lemma \ref{lem:Exp-O2NC:criterion-reduction} implies that, with $c=O(\epsilon^{-1})$, the algorithm automatically achieves the optimal rate of $O(\epsilon^{-4})$ for smooth objectives \cite{arjevani2019lower}. Similarly, with $c=O(1)$, it achieves the optimal rate of $O(\epsilon^{-7/2})$ for second-order smooth objectives \cite{arjevani2020second}. 

Secondly, we show in the following lemma that $(c,\epsilon)$-stationary points can also be reduced to Goldstein stationary points when the objective is Lipschitz.

\begin{restatable}{lemma}{GoldsteinReduction}
\label{lem:Exp-O2NC:goldstein-reduction}
Suppose $F$ is $G$-Lipschitz. For any $c,\epsilon,\delta>0$, a $(c,\epsilon)$-stationary point is also a $(\delta,\epsilon')$-Goldstein stationary point where $\epsilon' = (1+\frac{2G}{c\delta^2})\epsilon$.
\end{restatable}



\subsection{Online-to-non-convex Conversion}

\begin{algorithm}[t]
\caption{O2NC \cite{cutkosky2023optimal}}
\label{alg:o2nc-original}
\begin{algorithmic}[1]
    \STATE \textbf{Input:} OCO algorithm $\calA$, initial state $\vecx_0$, parameters $N,K,T\in\NN$ such that $N=KT$.
    \FOR {$n\gets 1,2,\ldots,N$}
        \STATE Receive $\Delta_n$ from $\calA$.
        \STATE Update $\vecx_n \gets \vecx_{n-1} + \Delta_n$ and $\vecw_n\gets\vecx_{n-1}+s_n\Delta_n$, where $s_n\sim\Unif([0,1])$ i.i.d.
        \STATE Compute $\vecg_n \gets \nabla f(\vecx_n, z_n)$.
        \STATE Send loss $\ell_n(\Delta) = \langle \vecg_n,\Delta\rangle$ to $\calA$.
        
        % \STATE 
        // For output only (update every $T$ iteration):
        \STATE If $n=kT$, compute $\overline\vecw_k = \frac{1}{T}\sum_{t=0}^{T-1} \vecw_{n-t}$.
    \ENDFOR
    % \State Compute $\vecx^k \gets \frac{1}{T} \sum_{n\in I_k} \vecx_n$ where $I_k=[(k-1)T+1,kT]$.
    \STATE Output $\overline \vecw \sim \Unif(\{\overline\vecw_k:k\in[K]\})$.
\end{algorithmic}
\end{algorithm}

Since our algorithm is an extension of the online-to-non-convex conversion (O2NC) technique proposed by \cite{cutkosky2023optimal}, we briefly review the original O2NC algorithm. The pseudocode is outlined in Algorithm \ref{alg:o2nc-original}, with minor adjustments in notations for consistency with our presentation. 

At its essence, O2NC shifts the challenge of optimizing a non-convex and non-smooth objective into minimizing regret. The intuition is as follows. By adding a uniform perturbation $s_n\in[0,1]$, $\langle \nabla f(\vecx_{n-1}+s_n\Delta_n,z_n), \Delta_n\rangle = \langle \vecg_n,\Delta_n\rangle$ is an unbiased estimator of $F(\vecx_n)-F(\vecx_{n-1})$, effectively capturing the ``training progress''. Consequently, by minimizing the regret, which is equivalent to minimizing $\sum_{n=1}^N\langle \vecg_n,\Delta_n\rangle$, the algorithm automatically identifies the most effective update step $\Delta_n$.


\subsection{Paper Organization}

In Section \ref{sec:alg}, we present the general online-to-non-convex framework, \alg. We  first explain the intuitions behind the algorithm design, and then we provide the convergence analysis in Section \ref{sec:alg-conv}. In Section \ref{sec:sgdm}, we provide an explicit instantiation of our framework, and see that the resulting algorithm is essentially the standard SGDM.
% In section \ref{sec:adam}, we further apply adaptive learning rate.
In Section \ref{sec:lower}, we present a lower bound for finding $(c,\epsilon)$-stationary point.
In Section \ref{sec:experiment}, we present empirical evaluations.